% TODO make hw stylesheet
\documentclass[11pt]{scrartcl}
\usepackage[a4paper, total={6in, 8in}]{geometry}
\usepackage[usenames,dvipsnames,svgnames]{xcolor}
\usepackage[shortlabels]{enumitem}
\usepackage[framemethod=TikZ]{mdframed}
\usepackage{amsmath,amssymb,amsthm}
\usepackage[colorlinks]{hyperref}
\usepackage{multicol}
\usepackage{tabularx}
\usepackage{stmaryrd}

\hypersetup{
	colorlinks=true,
	linkcolor=blue,
	filecolor=magenta,
	urlcolor=magenta,
	pdftitle={MAT 534 HW}
}

\usepackage{microtype}
\usepackage{mathtools}
\usepackage[headsepline]{scrlayer-scrpage}
\usepackage{thmtools}
\usepackage[textsize=footnotesize]{todonotes}

\mdfdefinestyle{mdbluebox}{roundcorner=10pt,innerbottommargin=9pt,
	linecolor=blue,backgroundcolor=TealBlue!5,}
\declaretheoremstyle[headfont=\sffamily\bfseries\color{MidnightBlue},
mdframed={style=mdbluebox},]{thmbluebox}

\mdfdefinestyle{mdbloobox}{frametitlefont=\bfseries,innerbottommargin=8pt,
	nobreak=true,backgroundcolor=TealBlue!5,linecolor=blue,}
\declaretheoremstyle[headfont=\bfseries\color{MidnightBlue},
mdframed={style=mdbloobox},headpunct={\\[3pt]},postheadspace=0pt,]{thmbloobox}

\mdfdefinestyle{mdredbox}{frametitlefont=\bfseries,innerbottommargin=8pt,
	nobreak=true,backgroundcolor=Salmon!5,linecolor=RawSienna,}
\declaretheoremstyle[headfont=\bfseries\color{RawSienna},
mdframed={style=mdredbox},headpunct={\\[3pt]},postheadspace=0pt,]{thmredbox}

\mdfdefinestyle{mdgreenbox}{linecolor=ForestGreen,backgroundcolor=ForestGreen!5,
	linewidth=2pt,rightline=false,leftline=true,topline=false,bottomline=false,}
\declaretheoremstyle[headfont=\bfseries\sffamily\color{ForestGreen!70!black},
mdframed={style=mdgreenbox},headpunct={ --- },]{thmgreenbox}

\mdfdefinestyle{mdorangebox}{linecolor=RedOrange,backgroundcolor=RedOrange!5,
	linewidth=2pt,rightline=false,leftline=true,topline=false,bottomline=false,}
\declaretheoremstyle[headfont=\bfseries\sffamily\color{RedOrange!70!black},
mdframed={style=mdorangebox},headpunct={ --- },]{thmorangebox}

\mdfdefinestyle{mdblackbox}{linecolor=black,backgroundcolor=RedViolet!5!gray!5,
	linewidth=3pt,rightline=false,leftline=true,topline=false,bottomline=false,}
\declaretheoremstyle[mdframed={style=mdblackbox}]{thmblackbox}

\declaretheoremstyle[spaceabove=6pt,spacebelow=6pt]{basehead}

\declaretheorem[style=basehead,name=Problem]{problem}
\declaretheorem[style=thmredbox,name=Problem,numbered=no]{problem*}
\declaretheorem[style=thmbloobox,name=Setup]{setup} % for config geo
\declaretheorem[style=thmredbox,name=Example,sibling=problem]{example}
\declaretheorem[style=thmredbox,name=$\bigstar$ Problem,sibling=problem]{niceprob}
\declaretheorem[style=thmbluebox,name=Theorem,sibling=problem]{theorem}
\declaretheorem[style=thmbluebox,name=Corollary,sibling=theorem]{corollary}
\declaretheorem[style=thmbluebox,name=Proposition,sibling=theorem]{prop}
\declaretheorem[style=thmbluebox,name=Lemma,numberwithin=problem]{lemma}
\declaretheorem[style=thmbluebox,name=Theorem,numbered=no]{theorem*}
\declaretheorem[style=thmbluebox,name=Lemma,numbered=no]{lemma*}
\declaretheorem[style=thmgreenbox,name=Claim,numberwithin=problem]{claim}
\declaretheorem[style=thmgreenbox,name=Claim,numbered=no]{claim*}
\declaretheorem[style=thmblackbox,name=Remark,numberwithin=problem]{remark}
\declaretheorem[style=thmblackbox,name=Remark,numbered=no]{remark*}
\declaretheorem[style=thmgreenbox,name=Definition,sibling=theorem]{definition}
\declaretheorem[style=thmgreenbox,name=Definition,numbered=no]{definition*}
\declaretheorem[style=thmorangebox,name=Warning,sibling=theorem]{warning}
\declaretheorem[style=thmorangebox,name=Warning,numbered=no]{warning*}
\declaretheorem[style=thmorangebox,name=Note,numbered=no]{note}
\declaretheorem[style=thmblackbox,name=Exercise,sibling=problem]{exercise}
\declaretheorem[style=thmblackbox,name=Exercise,numbered=no]{exercise*}
\declaretheorem[style=thmblackbox,name=Question,sibling=problem]{question}
\declaretheorem[style=thmblackbox,name=Question,numbered=no]{question*}

\addtolength{\textheight}{3.14cm}
\providecommand{\ol}{\overline}
\providecommand{\ceil}[1]{\left\lceil #1 \right\rceil}
\providecommand{\floor}[1]{\left\lfloor #1 \right\rfloor}
\newcommand{\dd}{\mathrm{d}}
\providecommand{\ul}{\underline}
\providecommand{\eps}{\varepsilon}
\providecommand{\half}{\frac{1}{2}}
\providecommand{\dang}{\measuredangle} %% Directed angle
\providecommand{\CC}{\mathbb C}
\providecommand{\FF}{\mathbb F}
\providecommand{\NN}{\mathbb N}
\providecommand{\QQ}{\mathbb Q}
\providecommand{\RR}{\mathbb R}
\providecommand{\ZZ}{\mathbb Z}
\providecommand{\dg}{^\circ}
\providecommand{\ii}{\item}
\providecommand{\setdiff}{\smallsetminus}
\providecommand{\alert}[1]{{\sffamily\textbf{\textcolor{blue}{#1}}}}
\providecommand{\inv}{^{-1}}
\providecommand{\mb}[1]{\mathbf{#1}}

\reversemarginpar

\newenvironment{soln}{\begin{proof}[Solution]}{\end{proof}}
\newenvironment{walkthrough}{\noindent \textbf{Walkthrough.}}{\medskip}
\newlist{walk}{enumerate}{3}
\setlist[walk]{label=\bfseries (\alph*)}

\providecommand{\pow}{\operatorname{Pow}}
\providecommand{\vocab}[1]{{\textbf{\textcolor{ForestGreen}{#1}}}}
\providecommand{\lcm}{\operatorname{lcm}}
\providecommand{\abs}[1]{\left|#1\right|}

\usepackage{asymptote}
\begin{asydef}
	defaultpen(fontsize(10pt));
	unitsize(1cm);
	size(8cm); // set a reasonable default
	usepackage("amsmath");
	usepackage("amssymb");
	settings.tex="pdflatex";
	settings.outformat="pdf";
	// Replacement for olympiad+cse5 which is not standard
	import geometry;
	// recalibrate fill and filldraw for conics
	void filldraw(picture pic = currentpicture, conic g, pen fillpen=defaultpen, pen drawpen=defaultpen)
	{ filldraw(pic, (path) g, fillpen, drawpen); }
	void fill(picture pic = currentpicture, conic g, pen p=defaultpen)
	{ filldraw(pic, (path) g, p); }
	// some geometry
	pair foot(pair P, pair A, pair B) { return foot(triangle(A,B,P).VC); }
	pair orthocenter(pair A, pair B, pair C) { return orthocentercenter(A,B,C); }
	pair centroid(pair A, pair B, pair C) { return (A+B+C)/3; }
	// cse5 abbreviations
	path CP(pair P, pair A) { return circle(P, abs(A-P)); }
	path CR(pair P, real r) { return circle(P, r); }
	pair IP(path p, path q) { return intersectionpoints(p,q)[0]; }
	pair OP(path p, path q) { return intersectionpoints(p,q)[1]; }
	path Line(pair A, pair B, real a=0.6, real b=a) { return (a*(A-B)+A)--(b*(B-A)+B); }
	// cse5 more useful functions
	picture CC() {
		picture p=rotate(0)*currentpicture;
		currentpicture.erase();
		return p;
	}
	pair MP(Label s, pair A, pair B = plain.S, pen p = defaultpen) {
		Label L = s;
		L.s = "$"+s.s+"$";
		label(L, A, B, p);
		return A;
	}
	pair Drawing(Label s = "", pair A, pair B = plain.S, pen p = defaultpen) {
		dot(MP(s, A, B, p), p);
		return A;
	}
	path Drawing(path g, pen p = defaultpen, arrowbar ar = None) {
		draw(g, p, ar);
		return g;
	}
\end{asydef}

\usepackage{mathrsfs}

\ihead{\footnotesize MAT 534 HW01}
\ohead{\footnotesize September 10, 2025}

\title{\vspace{-2em}MAT 534 HW01}
\author{\large Aditya Pahuja}
\date{\large Due 09/10}
\begin{document}
\maketitle

\begin{problem}
    Suppose that $G$ is a set with an associative operation,
    written as multiplication. Show that $G$ is a group if and only if,
    for every $a$ and $b$ in $G$, the equations
    $xa = b$ and $ay = b$ have solutions $x,y\in G$.
\end{problem}

\begin{soln}
    If $G$ is a group, then $xa = b$ has the solution $x = ba\inv$
    and $ay = b$ has the solution $y = a\inv b$.

    Conversely suppose that $xa = b$ and $ay = b$
    have solutions for all $a,b\in G$.
    Fixing $a\in G$, let $e$ be an element of $G$ such that $ea = a$.
    Given $b\in G$, pick $y$ so that $ay = b$. Then,
    \begin{equation*}
	eb = e(ay) = (ea)y = ay = b
    \end{equation*}
    so multiplication on the left by $e$ fixes all elements of $G$.
    Similarly, there exists an $e'\in G$ such that $ae' = a$,
    and then, given an arbitrary $b$ in $G$, if $xa = b$, we get
    \begin{equation*}
        be' = (xa)e' = x(ae') = xa = b
    \end{equation*}
    so multiplication on the right by $e'$ fixes all elements of $G$.
    This implies that
    \begin{equation*}
        e = ee' = e'
    \end{equation*}
    so in fact $e$ is the identity element for the operation on $G$.

    It remains to check the existence of inverses in $G$.
    For some $a\in G$, we are given that there exists a left inverse
    $x$ (so $xa = e$) and a right inverse $y$ (so $ay = e$);
    we need to check that $x = y$. This follows from
    \begin{equation*}
        x = x(ay) = (xa)y = y,
    \end{equation*}
    so for any $a\in G$ there exists $x$ such that $ax = xa = e$.
\end{soln}

\begin{problem}
    Suppose $G$ is a group in which $x^2 = 1$ for every $x\in G$.
    Show that $G$ is abelian.
\end{problem}

\begin{soln}
    Since $x^2 = 1$ for all $x$, each element of $G$ is its own inverse.
    This means that for any $a,b\in G$,
    \begin{equation*}
	ab = (ab)\inv = b\inv a\inv = ba
    \end{equation*}
    so the operation of $G$ is commutative.
\end{soln}

\pagebreak

\begin{problem}
    Let $G$ be a group and $o$ be any element of $G$.
    Define a new operation $*$ on $G$ by the formula $x*y = xoy$.
    \begin{enumerate}[(a)]
        \ii Show that $G$ is a group under the new operation $*$.
	What is the identity of $G$ for the new operation?
	If $x\in G$, what is the inverse of $x$ for the new operation?
	\ii If $G$ is a group and $o\in G$,
	let $G_o$ denote the group defined in (a).
	Show that $G\cong G_o$.
    \end{enumerate}
\end{problem}

\begin{soln}
    (a) First, $*$ is associative because
    \begin{equation*}
	(x*y)*z = (xoy)oz = xo(yoz) = x*(y*z),
    \end{equation*}
    using the fact that the operation of $G$ is of course associative.
    Second, the identity element for the operation $*$ is $o\inv$,
    since for any $x\in G$,
    \begin{equation*}
        x*o\inv = xoo\inv = x,\quad o\inv*x = o\inv ox = x.
    \end{equation*}
    Finally, given $x\in G$, the inverse of $x$ for $*$ is
    $(oxo)\inv = o\inv x\inv o\inv$, as it is easy to see that
    \begin{equation*}
        xo(oxo)\inv = (oxo)\inv xo = o\inv.
    \end{equation*}
    Thus $(G,*)$ satisfies the axioms of a group.

    (b) Define $\varphi\colon G\to G_o$ by $\varphi(x) = o\inv x$
    (where the inverse $o\inv$ is with respect to the operation of $G$).
    This is a bijective function because it is invertible via the map
    $y\mapsto oy$, and it is a homomorphism because
    \begin{equation*}
        \varphi(xy) = o\inv (xy) = (o\inv x)o(o\inv y) =
	\varphi(x)*\varphi(y),
    \end{equation*}
    so $\varphi$ is an isomorphism $G\to G_o$.
\end{soln}

\begin{problem}
    \begin{enumerate}[(a)]
        \ii Let $G$ be the set of linear polynomials $f(x) = ax+b$,
	where $a,b\in\RR$ and $a\ne0$. Show that $G$ is a group
	under composition of functions.
	\ii Let
	\begin{equation*}
	    H = \left\{ \begin{pmatrix}
		    a&b\\0&1
	    \end{pmatrix} : a,b\in\RR,\, a\ne 0\right\}.
	\end{equation*}
	Show that $H$ is a group under matrix multiplication.
	\ii Show that $G$ and $H$ are isomorphic.
    \end{enumerate}
\end{problem}

\begin{soln}
    (a) Composing two linear polynomials gives another linear polynomial:
    if $f(x) = ax+b$ and $g(x) = cx+d$, then
    \begin{equation*}
	(g\circ f)(x) = c(ax + b) + d = acx + (bc+d),
    \end{equation*}
    so $G$ is closed under composition.
    Function composition is of course associative,
    with identity being the identity function $e(x) = x$ for all $x$.
    Finally, the inverse of $f(x) = ax + b$ is
    \begin{equation*}
	f\inv(x) = \frac xa - \frac ba,
    \end{equation*}
    since
    \begin{align*}
	(f\inv\circ f)(x) &= \frac{ax+b}a - \frac ba = x + \frac ba - \frac ba = x \\
	(f\circ f\inv)(x) &= a\left(\frac xa - \frac ba\right) + b
	= x - b + b = x.
    \end{align*}

    (b) To see that $H$ is closed under matrix multiplication, we compute
    \begin{equation*}
        \begin{pmatrix}
	    c&d\\ 0&1
        \end{pmatrix}
        \begin{pmatrix}
	    a&b\\ 0&1
        \end{pmatrix} =
        \begin{pmatrix}
	    ac&bc+d\\ 0&1
        \end{pmatrix}.
    \end{equation*}
    Matrix multiplication is of course associative,
    with identity given by the identity matrix.
    Finally,
    \begin{equation*}
        \begin{pmatrix}
	    a&b\\0&1
	\end{pmatrix}\inv =
	\begin{pmatrix}
	    1/a & -b/a\\0&1
	\end{pmatrix},
    \end{equation*}
    which can be verified easily.
    (Less directly, we can also just note that the inverse in
    $GL_2(\RR)$ of an element of $G$ has $(0\; 1)$ as its bottom row
    due to how the row reduction algorithm works.)

    (c) $G$ and $H$ are isomorphic with isomorphism $\varphi\colon G\to H$ given by
    \begin{equation*}
        \varphi(x\mapsto ax+b) = \begin{pmatrix}
	    a&b\\0&1
        \end{pmatrix}.
    \end{equation*}
    This is clearly invertible hence bijective, as the first row of the matrix
    exactly encodes the coefficients of the corresponding linear polynomial in $G$.
    To check that this is a homomorphism, if $f(x) = ax+b$
    and $g(x) = cx+d$ with $a,c\ne0$, then
    \begin{equation*}
        \varphi(g)\varphi(f) = \begin{pmatrix}
	    c&d\\0&1
        \end{pmatrix}
	\begin{pmatrix}
	    a&b\\0&1
	\end{pmatrix} =
	\begin{pmatrix}
	    ac & bc+d \\ 0&1
	\end{pmatrix} =
	\varphi(g\circ f),
    \end{equation*}
    recalling from part (a) that $(g\circ f)(x) = acx + (bc + d)$.
\end{soln}

\begin{problem}
    Let $\CC^\times$ be the multiplicative group of the complex numbers,
    i.e. the non-zero elements of $\CC$, under multiplication.
    Show that $\CC^\times$ is isomorphic to the direct product
    of the two additive groups $\RR$ and $\RR/\ZZ$.
\end{problem}

\begin{soln}
    Consider the map
    \begin{align*}
        \varphi\colon \RR\times\RR &\to \CC^\times \\
	(r,\theta) &\mapsto e^{r + 2\pi i\theta} =
	e^r(\cos(2\pi\theta) + i\sin(2\pi\theta)).
    \end{align*}
    This is clearly a homomorphism, as
    \begin{equation*}
	\varphi(r+s,\theta+\psi) = e^{r+s + 2\pi i(\theta + \psi)}
	= e^{r + 2\pi i\theta} e^{s+2\pi i\theta}
	= \varphi(r,\theta)\varphi(s,\psi).
    \end{equation*}
    Furthermore, the kernel of this homomorphism is the set of $(r,\theta)$ such that
    \begin{equation*}
	e^{r + 2\pi\theta} = 1 \iff e^r = 1,\; 2\pi\theta\in 2\pi\ZZ
    \end{equation*}
    which is equivalent to $\theta\in\ZZ$ and $r = 0$, so
    \begin{equation*}
	\CC^\times\cong\frac{\RR\times\RR}{\langle (0,1)\rangle}.
    \end{equation*}
    In turn, by the homomorphism $\RR\times\RR\to\RR\times(\RR/\ZZ)$
    defined by $(a,b)\mapsto (a,b\bmod 1)$, whose kernel is
    $\langle (0,1)\rangle$, we get
    \begin{equation*}
	\CC^\times\cong\frac{\RR\times\RR}{\langle (0,1)\rangle}
	\cong \RR\times(\RR/\ZZ),
    \end{equation*}
    which is what we wanted to show.
\end{soln}

\pagebreak

\begin{problem}
    Let $n\ge 2$ be an integer. Determine the center
    of the dihedral group $D_{2n}$.
\end{problem}

\begin{soln}
    Recall that
    \begin{equation*}
        D_{2n} = \langle r,d \mid r^n = d^2 = 1,\, rd = dr\inv\rangle
	= \{1, r,\dots, r^{n-1}, d, rd, \dots, r^{n-1}d\}.
    \end{equation*}
    When $n = 2$, it is easy to check that the group is abelian,
    hence $Z(D_{2n}) = D_{2n}$. Assume henceforth that $n > 2$.

    First, note that $r^kd$ is always not in the center of $D_{2n}$,
    seeing as $r\cdot r^kd = r^kd\cdot r\inv \ne r^kd\cdot r$ because $r\ne r\inv$
    (since $r^2=1$ would imply $n=2$).

    Thus $Z(D_{2n})$ is a subgroup of $\{1, r, \dots, r^{n-1}\}$.
    Obviously the powers of $r$ all commute with each other,
    so we need to check which powers of $r$ commute with the $r^kd$.
    Take $0 \le m < n$. Then
    \begin{equation*}
	r^kd\cdot r^m = r^m\cdot r^kd = r^kd\cdot r^{-m}
    \end{equation*}
    if and only if $r^{2m} = 1$. This requires $n\mid 2m$,
    which, along with the bound $0\le 2m < 2n$, implies
    $2m = 0$ or $2m = n$. The former case tells us that $1$ is in the center (shocking),
    while the latter tells us that, if $n$ is even, then $r^{n/2}$ is in the center.
    Therefore
    \begin{equation*}
	Z(D_{2n}) = \begin{cases}
	    D_4 & n = 2\\
	    \{1\} & n\text{ is odd}\\
	    \{1, r^{n/2}\} & n\text{ is even}.
	\end{cases}
    \end{equation*}
\end{soln}

\begin{problem}
    Let $H$ and $K$ be subgroups of $G$, with $K$ normal in $G$.
    \begin{enumerate}[(a)]
	\ii Show that $HK = \{hk : h\in H,\, k\in G\}$
	is a subgroup of $G$.
	\ii Suppose that $H$ is also normal.
	Show that $HK$ is a normal subgroup of $G$.
    \end{enumerate}
\end{problem}

\begin{soln}
    (a) To show that $HK$ is a subgroup of $G$,
    we need to show that if $g$ and $g'$ are in $HK$,
    then so is $g(g')\inv$.
    In this spirit, let $h_1k_1$ and $h_2k_2$ be elements of $HK$,
    with $h_1,h_2\in H$ and $k_1,k_2\in K$. Then,
    \begin{equation*}
        h_1k_1(h_2k_2)\inv = h_1k_1k_2\inv h_2\inv =
	h_1 h_2\inv (h_2 k_1k_2\inv h_2\inv).
    \end{equation*}
    Since $K$ is normal in $G$, $k_1k_2\inv\in K$ implies that the conjugate
    $h_2 k_1k_2\inv h_2\inv$ is also in $K$,
    so we have expressed $h_1k_1(h_2k_2)\inv$ as the product of
    an element $h_1h_2\inv$ of $H$ and an element
    $h_2k_1k_2\inv h_2\inv$ of $K$, hence it is in $HK$ as desired.

    (b) If $H$ is also normal, then for any $hk\in HK$
    with $h\in H$, $k\in K$ and any $g\in G$, we get
    \begin{equation*}
	g hk g\inv = \underbrace{(ghg\inv)}_{\in H}
	\underbrace{(gkg\inv)}_{\in K}\in HK,
    \end{equation*}
    so $HK$ is closed under conjugation, hence normal.
\end{soln}

\pagebreak

\begin{problem}
    The \vocab{commutator subgroup} of a group $G$ is
    the group $G'$ generated by the elements $[x,y] = x\inv y\inv xy$,
    where $x$ and $y$ vary through the elements of $G$.
    \begin{enumerate}[(a)]
        \ii Show that $G'$ is normal in $G$.
	\ii Show that $G/G'$ is an abelian group.
	\ii Let $H$ be a subgroup of $G$ containing $G'$.
	Show that $H$ is normal in $G$, and that $G/H$ is abelian.
	\ii Suppose that $N$ is a normal subgroup of $G$
	such that $G/N$ is abelian. Show that $G'\subseteq N$.
    \end{enumerate}
\end{problem}

\begin{soln}
    (a) First, observe that the conjugate $h\inv[x,y]h$ of a commutator
    is also a commutator $[h\inv xh, h\inv yh]$.
    Then, an element $g$ of $G'$ is a product of finitely many commutators, say
    \begin{equation*}
	g = [x_1,y_1][x_2,y_2]\cdots[x_n,y_n].
    \end{equation*}
    For any $h\in G$,
    \begin{equation*}
	h\inv gh = (h\inv [x_1,y_1]h)(h\inv[x_2,y_2]h)\cdots(h\inv[x_n,y_n]h)
	= \prod_{i=1}^n [h\inv x_ih, h\inv y_ih]
    \end{equation*}
    is thus a product of commutators, so $G'$ is closed under conjugation,
    hence is normal.

    (b) Let $\varphi\colon G\to G/G'$ be the projection homomorphism,
    i.e. $\varphi(x)$ gets sent to the coset of $G'$ containing $x$;
    this is a surjective homomorphism with kernel $G'$.
    Then, given $g,h\in G/G'$, let $x$ and $y$ be elements of
    the cosets $g$, $h$; we have
    \begin{equation*}
        gh = \varphi(x)\varphi(y) = \varphi(xy)
	= \varphi([x\inv, y\inv]yx) = \varphi([x\inv, y\inv])\varphi(xy)
	= \varphi(xy) = hg
    \end{equation*}
    because the commutator $[x\inv, y\inv]$ is in $G'$, so $G/G'$ is abelian.

    (c) If $h\in H$, then $h\inv\in H$ as well. Thus for any $g\in G$,
    \begin{equation*}
	g\inv hg\in H \iff h\inv (g\inv hg) = [h,g]\in H,
    \end{equation*}
    but we know that $[h,g]\in H$ because $H$ contains all the commutators in $G$.
    Thus, $H$ is closed under conjugation, hence normal.

    It is straightforward to show that $G/H$ is abelian,
    with the exact same logic as part (b).
    Let $\varphi\colon G\to G/H$ be the projection homomorphism,
    Then, given $g,h\in G/H$, let $x$ and $y$ be elements of
    the cosets $g$, $h$; we have
    \begin{equation*}
        gh = \varphi(x)\varphi(y) = \varphi(xy)
	= \varphi([x\inv, y\inv]yx) = \varphi([x\inv, y\inv])\varphi(xy)
	= \varphi(xy) = hg
    \end{equation*}
    because the commutator $[x\inv, y\inv]$ is in $G'\subseteq H$,
    so $G/H$ is abelian.

    (d) If $G/N$ is abelian, let $\varphi\colon G\to G/N$ be
    the projection homomorphism. Then, for any $x,y\in G$, we have
    \begin{equation*}
        \varphi(x)\varphi(y) = \varphi(y)\varphi(x)
	\iff \varphi(x\inv y\inv xy) = 1,
    \end{equation*}
    so every commutator is in $N$, hence the subgroup generated by the commutators,
    $G'$, is contained in $N$.
\end{soln}


















\end{document}
